\documentclass[11pt,letterpaper]{article}
\usepackage[T1]{fontenc}
\usepackage[utf8]{inputenc}
\usepackage[margin=0.88in,top=0.84in,bottom=0.83in,headheight=13pt,headsep=20pt,footskip=28pt]{geometry}
\usepackage{newpxtext}
\usepackage{amsmath,amsthm}
\usepackage{newpxmath}
% newpxtext selects TeX Gyre Heros (qhv), scaled to .94, for sans serif.
\usepackage{microtype}
\usepackage{xcolor}
\definecolor{Forest}{HTML}{30392C}
\definecolor{Olive}{HTML}{687144}
\definecolor{Muted}{HTML}{65685F}
\definecolor{Sage}{HTML}{CBD0BC}
\definecolor{Pale}{HTML}{F2F4EB}
\usepackage{mathtools}
\usepackage{booktabs,tabularx,longtable,array}
\usepackage{enumitem}
\usepackage{titlesec}
\usepackage{fancyhdr}
\usepackage[most]{tcolorbox}
\usepackage{needspace}
\usepackage{xurl}
\usepackage[colorlinks=true,linkcolor=Olive,citecolor=Olive,urlcolor=Olive,bookmarksnumbered=true,pdfencoding=auto,psdextra]{hyperref}
\hypersetup{pdftitle={Cover exactingness, definability, and conditional inconsistency},pdfauthor={Prepared for Vladimir Reshetnikov},pdfsubject={Proof-oriented follow-up: HCD regularity, covering obstructions, grounds, and an I0 equiconsistency corollary},pdfkeywords={cover exacting, HCD, strongly compact, Ground Axiom, ultraexacting, equiconsistency}}
\urlstyle{same}
\setlength{\parindent}{0pt}
\setlength{\parskip}{5.1pt plus 1pt minus .4pt}
\linespread{1.035}
\setlength{\emergencystretch}{2em}
\setcounter{tocdepth}{1}
\setcounter{secnumdepth}{2}
\titleformat{\section}{\Large\bfseries\color{Forest}}{\thesection}{.6em}{}
\titleformat{\subsection}{\large\bfseries\color{Forest}}{\thesubsection}{.55em}{}
\titleformat{\subsubsection}{\normalsize\bfseries\color{Forest}}{}{0pt}{}
\titlespacing*{\section}{0pt}{18pt plus 3pt minus 2pt}{8pt}
\titlespacing*{\subsection}{0pt}{12pt plus 2pt minus 1pt}{5pt}
\titlespacing*{\subsubsection}{0pt}{9pt}{3pt}
\setlist[itemize]{leftmargin=1.4em,itemsep=2pt,topsep=3pt,parsep=0pt}
\setlist[enumerate]{leftmargin=1.7em,itemsep=3pt,topsep=4pt,parsep=0pt}
\pagestyle{fancy}
\fancyhf{}
\fancyhead[L]{\sffamily\fontsize{9}{11}\selectfont\color{Muted}LARGE CARDINALS / DEFINABILITY BOUNDARIES}
\fancyhead[R]{\sffamily\fontsize{9}{11}\selectfont\color{Muted}18 SEP 2026}
\fancyfoot[L]{\small\color{Muted}Definability and consistency boundaries}
\fancyfoot[R]{\small\color{Muted}\thepage}
\renewcommand{\headrulewidth}{.35pt}
\renewcommand{\footrulewidth}{0pt}
\fancypagestyle{plain}{\fancyhf{}\fancyfoot[L]{\small\color{Muted}Definability and consistency boundaries}\fancyfoot[R]{\small\color{Muted}\thepage}\renewcommand{\headrulewidth}{0pt}}
\tcbset{colback=Pale,colframe=Sage,colbacktitle=Sage,coltitle=Forest,fonttitle=\bfseries,boxrule=.5pt,arc=3pt,left=9pt,right=9pt,top=6pt,bottom=6pt,toptitle=3pt,bottomtitle=3pt,before skip=8pt,after skip=8pt,breakable}
\newtcolorbox{assessment}[1]{title={#1}}
\theoremstyle{definition}
\newtheorem{definition}{Definition}[section]
\newtheorem{theorem}[definition]{Theorem}
\newtheorem{proposition}[definition]{Proposition}
\newtheorem{lemma}[definition]{Lemma}
\newtheorem{corollary}[definition]{Corollary}
\newcommand{\ZFC}{\mathsf{ZFC}}
\newcommand{\ZF}{\mathsf{ZF}}
\newcommand{\AC}{\mathsf{AC}}
\newcommand{\DC}{\mathsf{DC}}
\newcommand{\HOD}{\mathrm{HOD}}
\newcommand{\OD}{\mathrm{OD}}
\newcommand{\HH}{\mathsf{HH}}
\newcommand{\Ex}{\operatorname{Ex}}
\newcommand{\UEx}{\operatorname{UEx}}
\newcommand{\Ext}{\operatorname{Ext}}
\newcommand{\SC}{\operatorname{SC}}
\newcommand{\CEx}{\operatorname{CEx}}
\newcommand{\Con}{\operatorname{Con}}
\newcommand{\crit}{\operatorname{crit}}
\newcommand{\cf}{\operatorname{cf}}
\newcommand{\ot}{\operatorname{ot}}
\newcommand{\Ord}{\mathrm{Ord}}
\newcommand{\Card}{\mathrm{Card}}
\newcommand{\rank}{\operatorname{rank}}
\newcommand{\id}{\mathrm{id}}
\newcommand{\restr}{\mathbin{\upharpoonright}}
\newcommand{\Izero}{\mathrm{I0}}
\newcommand{\Ione}{\mathrm{I1}}
\newcommand{\Itwo}{\mathrm{I2}}
\newcommand{\Ithree}{\mathrm{I3}}
\newcommand{\Isharp}{\mathrm{I0}^{\#}}
\newcommand{\Status}[1]{\textbf{Status.} #1}
\newcommand{\Priority}[1]{\textbf{Research priority.} #1}
\newcolumntype{Y}{>{\raggedright\arraybackslash}X}
\newcolumntype{P}[1]{>{\raggedright\arraybackslash}p{#1}}
\newcommand{\arxiv}[1]{\href{https://arxiv.org/abs/#1}{arXiv:#1}}
\newcommand{\doi}[1]{\href{https://doi.org/#1}{doi:\nolinkurl{#1}}}

\newcommand{\CD}{\mathrm{CD}}
\newcommand{\HCD}{\mathrm{HCD}}
\newcommand{\UF}{\mathrm{UF}}
\newcommand{\GA}{\mathsf{GA}}
\newcommand{\REx}{\operatorname{REx}}
\newcommand{\Cov}{\operatorname{Cov}}
\newcommand{\tc}{\operatorname{tc}}
\newcommand{\ran}{\operatorname{ran}}
\newcommand{\Pow}{\mathcal P}
\newcommand{\Bdry}{\mathsf{B}}
\newcommand{\Apx}{\mathsf{Apx}}
\newcommand{\Known}[1]{\textbf{Imported result.} #1}
\newtheorem{remark}[definition]{Remark}
\newtheorem{fact}[definition]{Fact}

\begin{document}
\thispagestyle{empty}
{\sffamily\bfseries\small\color{Olive}SET THEORY\quad /\quad RESEARCH FOLLOW-UP}

\vspace{13pt}
{\fontsize{28}{31}\selectfont\bfseries\color{Forest}Cover exactingness,\newline definability, and\newline conditional inconsistency\par}
\vspace{10pt}
{\fontsize{14}{18}\selectfont A cofinality barrier in the HCD hierarchy,\newline local covering obstructions, and an I0 comparison\par}
\vspace{14pt}
{\color{Muted}Prepared for Vladimir\hfill 18 September 2026\par}
\vspace{3pt}
{\color{Sage}\hrule height .6pt}
\vspace{12pt}

\textbf{Main result.} A $\gamma$-cover-exacting cardinal $\lambda$ is regular in
$\HCD(\gamma^+)$, the inner model of sets hereditarily definable from
$\gamma^+$-complete ultrafilters on ordinals. This yields two different
incompatibility results, depending on the position of a strongly compact cardinal.

\begin{assessment}{Below the cover-exacting cardinal: a forced gap}
If $\delta<\lambda$ is strongly compact and $\lambda$ is $\gamma$-cover exacting,
there is an unbounded $a\subseteq\lambda$ such that
\[
 a\in\HCD(\delta),\qquad \ot(a)<\delta,\qquad
 a\notin\HCD(\gamma^+).
\]
In fact, no member of $\HCD(\gamma^+)$ of ambient size less than $\lambda$
can contain $a$. Thus even a weak, explicitly specified local covering principle
between these two canonical inner models is inconsistent with the configuration.
\end{assessment}

\begin{assessment}{Above the cover bound: a proper ground is unavoidable}
If $\kappa>\gamma$ is strongly compact, then $\HCD(\kappa)$ is a proper ground
of $V$ in which $\lambda$ is regular. Consequently,
\[
 \ZFC+\GA+\exists\lambda\leq\gamma<\kappa\,
       [\CEx_\gamma(\lambda)\wedge\SC(\kappa)]
\]
is inconsistent. Here $\GA$ is the Ground Axiom. This is not a prohibition of
cover exactingness alone or of its coexistence with a larger strongly compact cardinal.
\end{assessment}

\textbf{Positive comparison.} An ultraexacting cardinal can consistently be the
first exacting cardinal, with every set of lower rank in HOD, while a single
countable cofinal set witnesses strong failure of approximation and covering
at the boundary. The full package is equiconsistent with $\Izero$.

\textbf{Attribution and scope.} The proofs below derive refinements from the
Blue--Goldberg cover-exacting framework and Goldberg's published HCD theorems.
The affirmative package is a corollary of the existing Aguilera--Bagaria--Goldberg--L\"ucke
forcing and equiconsistency results, not a new forcing construction. The bounded
HCD obstruction is the principal proposed refinement; historical priority is
not certified. The unrestricted strongly-compact-below-cover-exacting problem
is not resolved here.

\clearpage
\tableofcontents
\vspace{10pt}
\begin{assessment}{Reading guide}
Sections~\ref{sec:critical}--\ref{sec:barrier} prove the main definability barrier.
Section~\ref{sec:below} gives the local inconsistency theorem above a strongly compact cardinal.
Section~\ref{sec:ground} gives the Ground Axiom incompatibility in the opposite ordering.
Section~\ref{sec:positive} proves the affirmative equiconsistency package.
Section~\ref{sec:status} and Appendix~\ref{app:audit} distinguish the conclusions
from the unresolved targets in the supplied report.
\end{assessment}

\section{The research target and the resulting theorems}\label{sec:overview}

The supplied report singled out a specific gap in the cover-exacting programme:
the incompatibility with a smaller extendible cardinal is known, whereas the
corresponding question for a smaller strongly compact cardinal remains open in
the cited notes. Its proposed next step was to locate a definability or covering
condition sufficient to transfer the negative argument. We pursue exactly that
step, rather than repeat the candidate survey. \cite{previous,cover}

The key distinction is between \emph{obtaining a short cofinal set in a canonical
inner model} and \emph{making that set definable from a parameter that a relative
exacting witness can absorb}. Strong compactness already supplies the first
ingredient in $\HCD(\delta)$. The second ingredient fails at the relevant higher
completeness level. Keeping the two levels separate produces a precise
conditional inconsistency theorem.

For an infinite cardinal $\eta$, $\CD(\eta)$ consists of sets ordinal definable
from finitely many $\eta$-complete ultrafilters on ordinals; $\HCD(\eta)$ is its
hereditary part. All these classes are computed in the same ambient universe.
Definitions and source conventions are given in Section~\ref{sec:setup}.

\subsection{The definability barrier}

The central statement, proved as Theorem~\ref{thm:barrier}, is
\begin{equation}\label{eq:mainbarrier}
 \CEx_\gamma(\lambda)\quad\Longrightarrow\quad
 \begin{cases}
 \text{no unbounded }a\subseteq\lambda\text{ of order type }<\lambda
     \text{ belongs to }\CD(\gamma^+),\\
 \cf^{\HCD(\gamma^+)}(\lambda)=\lambda.
 \end{cases}
\end{equation}
The successor on $\gamma$ has a concrete role: a set of size at most $\gamma$
can impose at most $\gamma$ many simultaneous ultrafilter constraints, and
$\gamma^+$-completeness guarantees that their intersection remains nonempty.
This observation is turned into a club of small elementary substructures on
which the ultrafilter becomes ordinal definable.

The statement is stronger than regularity in HOD. It rules out every short
cofinal set definable from the specified complete-ultrafilter parameters,
not merely every ordinal-definable one. No Axiom of Choice inside
$\HCD(\gamma^+)$ is assumed.

\subsection{Two conditional negative conclusions}

For a strongly compact $\delta<\lambda$, put
\[
 N_0=\HCD(\delta),\qquad N_1=\HCD(\gamma^+).
\]
Then $N_1\subseteq N_0$ and Theorem~\ref{thm:gap} gives
\begin{equation}\label{eq:cofgap}
 \cf^{N_0}(\lambda)<\delta<\lambda=\cf^{N_1}(\lambda).
\end{equation}
More precisely, there is an $N_0$-set $a\subseteq\lambda$ of order type less than
$\delta$ that has no $N_1$-superset of ambient size less than $\lambda$.
Consequently, either local small-subset agreement or a much weaker local
covering principle between $N_1$ and $N_0$ is enough to contradict the
large-cardinal configuration. These principles are stated as additional
hypotheses, not inferred from strong compactness.

In the other ordering, a strongly compact $\kappa>\gamma$ provides a ground
$\HCD(\kappa)$ contained in $\HCD(\gamma^+)$. Equation~\eqref{eq:mainbarrier}
therefore makes $\lambda$ regular in that ground, although it has cofinality
$\omega$ in $V$. The ground must be proper. This proves the Ground Axiom
incompatibility in Corollary~\ref{cor:ga}.

\subsection{An affirmative consistency boundary}

Theorem~\ref{thm:eq} gives an explicit theory $\Bdry$ with
\[
 \Con(\ZFC+\Izero)\quad\longleftrightarrow\quad\Con(\ZFC+\Bdry).
\]
In a model of $\Bdry$, an ultraexacting $\lambda$ is the first exacting cardinal,
$V_\lambda\subseteq\HOD$, and a countable cofinal set $c\subseteq\lambda$ is
outside HOD although all its intersections with HOD-sets of HOD-size less than
$\lambda$ are finite. Thus very extensive rank agreement is compatible with
failure of precisely the covering and approximation principles that would
internalize the cofinal sequence.

This comparison supplies a control on the negative arguments. Rank agreement
below $\lambda$ is not a substitute for covering at $\lambda$, and an
inconsistency proof cannot import such covering without a separate argument.

\section{Definitions, formal setting, and imported results}\label{sec:setup}

\subsection{Ambient and internal conventions}

We work in $\ZFC$. All unadorned cardinalities, successor cardinals, clubs, and
stationary sets are ambient ones. For a transitive inner model $N$, internal
cardinality and cofinality are written $|x|^N$ and $\cf^N(\alpha)$.
If $a$ is a set of ordinals belonging to two transitive models, its increasing
order type $\ot(a)$ is absolute between them. In particular, the assertion
$\ot(a)<\lambda$ does not change when we move between the models below.

For an uncountable regular $\eta$ and a set $A$, let
\[
 \Pow_\eta(A)=\{\sigma\subseteq A:|\sigma|<\eta\}.
\]
A club on this set is unbounded under inclusion and closed under increasing
unions of nonzero length less than $\eta$. A stationary subset meets every
such club. In particular, for a fixed set structure in a countable language,
the elementary substructures of size less than $\eta$ containing a specified
finite list of elements form a club: take Skolem hulls for unboundedness and
use the elementary-chain theorem for closure.

The main cover-exacting arguments involve only set-sized embeddings, so their
elementarity can be expressed using satisfaction for set structures. The
$\Izero$ consistency comparison is used in the standard formalization of its
source. No elementarity for formulas containing an arbitrary embedding
predicate is silently added.

\subsection{Relative exactingness and the uniform cover bound}

\begin{definition}[Relative exacting witness]\label{def:rex}
For an ordinal $\alpha>\lambda$, a $\lambda$-exacting witness relative
to a set predicate $Y\subseteq V_\alpha$ is an elementary map
\[
 j:(X,X\cap Y)\longrightarrow(V_\alpha,Y)
\]
such that $(X,X\cap Y)\prec(V_\alpha,Y)$,
$V_\lambda\cup\{\lambda\}\subseteq X$, $\crit(j)<\lambda$, and
$j(\lambda)=\lambda$. Write $\REx(\lambda,Y)$ when such witnesses exist at
every sufficiently large height containing $Y$ as a set.
\end{definition}

This eventual-height convention is sufficient for every proof here. The
relative property in the source notes, which requires all applicable heights,
implies it. The predicate is not an arbitrary new class; it is a set and can
be regarded as an element at a larger rank. The domain $X$ need not be
transitive.

\begin{definition}[Cover exactingness]\label{def:cover}
A cardinal $\lambda$ is $\gamma$-cover exacting, written $\CEx_\gamma(\lambda)$,
if for every $\alpha>\lambda$ and $y\in V_\alpha$ there are
$Y\subseteq V_\alpha$ and a relative witness as in Definition~\ref{def:rex}
at that height, with
\[
 y\in Y,\qquad |Y|\leq\gamma.
\]
The cardinal $\gamma$ is fixed outside the height and parameter quantifiers.
\end{definition}

We use the source's definition, including its low-height convention; only high
instances are used below. We restrict attention to infinite cardinal bounds
$\gamma\geq\lambda$. It is important that $y\in Y$ is membership, not inclusion,
and that the witness need not fix the particular element $y$.
\cite[Definitions 2.6 and 3.1]{cover}

\begin{fact}[Stationary relative-witness characterization]\label{fact:stationary}
If $\CEx_\gamma(\lambda)$, then, for every sufficiently large $\alpha$, the set
\[
 \mathcal S_\alpha=
 \{\sigma\in\Pow_{\gamma^+}(V_\alpha):\REx(\lambda,\sigma)\}
\]
is stationary in $\Pow_{\gamma^+}(V_\alpha)$.
\end{fact}

\Known{This is the implication (1)$\Rightarrow$(3) of Blue--Goldberg's
Proposition 3.4, with the weaker eventual-height convention above.
\cite[Proposition 3.4]{cover} The stationary characterization is a substantive
input. It is not a consequence of choosing one arbitrary cover $Y$.}

Ordinary exactingness is obtained by omitting the predicate and cover clauses.
Ultraexactingness additionally requires the witness restriction
$j\restr V_\lambda$ to be an element of $X$. We use the equivalent cardinal-height
formulations of these two properties from the supplied report and
\cite{abl,abgl}.

\subsection{Complete-ultrafilter definability}

\begin{definition}\label{def:hcd}
An ultrafilter is $\eta$-complete if intersections of fewer than $\eta$ of its
members belong to it. All ultrafilters in this paper are proper.
Let $\UF_\eta(\Ord)$ be the class of $\eta$-complete ultrafilters on ordinal
underlying sets, and put
\[
 \CD(\eta)=\OD_{\UF_\eta(\Ord)},\qquad
 \HCD(\eta)=\HOD_{\UF_\eta(\Ord)}.
\]
Here the subscript allows finitely many \emph{members} of the indicated class
as parameters, together with ordinal parameters. It does not allow arbitrary
extra predicates. Let $\HCD=\bigcap_\eta\HCD(\eta)$.
\end{definition}

Equivalently, a single $\eta$-complete ultrafilter parameter suffices: finite
ultrafilter parameters and ordinal parameters can be combined using Fubini
products and ordinal codings. We do not need this reduction, since the club
argument handles finitely many parameters directly. Goldberg proves that
$\HCD(\eta)$ is an inner model of $\ZF$. \cite[Propositions 4.2--4.3]{unique}

Two elementary facts will be used repeatedly:
\begin{equation}\label{eq:monotone}
 \eta\leq\xi\quad\Longrightarrow\quad
 \CD(\xi)\subseteq\CD(\eta),\qquad
 \HCD(\xi)\subseteq\HCD(\eta);
\end{equation}
and, for every set of ordinals $a$,
\begin{equation}\label{eq:ordhereditary}
 a\in\CD(\eta)\quad\Longleftrightarrow\quad a\in\HCD(\eta).
\end{equation}
For the second equivalence, once $a$ itself is completely definable at that
level, everything else in its transitive closure is an ordinal and is already
ordinal definable. Monotonicity follows directly from ultrafilter completeness.

\subsection{The published HCD inputs}

\begin{fact}[Goldberg]\label{fact:hcd}
If $\delta$ is strongly compact, then $\HCD(\delta)$ is a model of $\ZFC$, is a
set-forcing ground of $V$, and has the $\delta$-cover property. In particular,
any ambient set of ordinals of size less than $\delta$ is contained in a set
of ordinals in $\HCD(\delta)$ of size less than $\delta$.
If $\delta$ is extendible, then $\HCD(\delta)=\HCD$.
\end{fact}

\Known{These are Theorems 4.4, 4.10, 4.14, and 4.15 of Goldberg's
\emph{The uniqueness of elementary embeddings}. \cite{unique}
The covering conclusion just stated is weaker than the full theorem and is
all that Section~\ref{sec:below} requires. The published ground and covering
theorems use \emph{strong compactness}; the supercompact hypothesis in the
shorter presentation of the lecture notes is not needed for these inputs.}

A \emph{ground} is a transitive inner model $N\models\ZFC$ such that
$V=N[G]$ for some set forcing $\mathbb P\in N$ and $N$-generic $G$.
The Ground Axiom $\GA$ says that there is no proper such ground.
This terminology only concerns set forcing, not arbitrary class forcing.

\section{Why a short cofinal predicate is fatal}\label{sec:critical}

This section makes explicit the elementary obstruction used later. The local
Kunen theorem is the only large-cardinal inconsistency theorem imported into
these proofs. The argument is the familiar critical-sequence obstruction,
with care taken over definability parameters. \cite{kunen,hkp,cover}

\begin{lemma}[The critical sequence reaches the cardinal height]\label{lem:critical}
Let $\lambda$ be a cardinal and let $e:V_\lambda\to V_\lambda$ be a nontrivial
elementary embedding with critical point $\kappa$. Set
\[
 \kappa_0=\kappa,\qquad \kappa_{n+1}=e(\kappa_n).
\]
Then $\sup_{n<\omega}\kappa_n=\lambda$. Moreover, no ordinal in
$[\kappa,\lambda)$ is fixed by $e$.
\end{lemma}

\begin{proof}
Let $\nu=\sup_n\kappa_n\leq\lambda$. Suppose $\nu<\lambda$. Because $\lambda$
is an infinite cardinal, the sequence and the ranks $V_{\nu+1},V_{\nu+2}$ lie
below the relevant height $\lambda$. The sequence belongs to $V_\lambda$,
so elementarity, $e(\omega)=\omega$, and the definition of its supremum give
$e(\nu)=\nu$. Restricting $e$ therefore yields a nontrivial elementary
self-embedding of $V_{\nu+2}$ with critical point below $\nu$. The local
Kunen theorem rules this out in $\ZFC$. Hence $\nu=\lambda$.

For $\kappa\leq\xi<\lambda$, choose $n$ with
$\kappa_n\leq\xi<\kappa_{n+1}$. Then
\[
 e(\xi)\geq e(\kappa_n)=\kappa_{n+1}>\xi.
\]
This proves the final assertion.
\end{proof}

In particular, such a cardinal height has ambient cofinality $\omega$.
A sufficiently high exacting witness restricts to an embedding of this kind:
$V_\lambda$ is definable from the fixed parameter $\lambda$, and every element
of $V_\lambda$ is in the source. Relativizing formulas to that set proves
elementarity of the restriction. This is not an appeal to transitivity of $X$.

\begin{lemma}[No fixed short cofinal set]\label{lem:fixedcofinal}
Let $j:X\to V_\alpha$ be a sufficiently high exacting witness at $\lambda$.
There is no unbounded $a\subseteq\lambda$ with $a\in X$, $j(a)=a$, and
$\ot(a)<\lambda$.
\end{lemma}

\begin{proof}
Let $\kappa=\crit(j)$ and $\rho=\ot(a)$. The order type and increasing
enumeration $f:\rho\to a$ are uniquely definable from $a$ in a sufficiently
high rank. Hence $\rho,f\in X$, $j(\rho)=\rho$, and $j(f)=f$.
Lemma~\ref{lem:critical} applied to $j\restr V_\lambda$ implies $\rho<\kappa$.
Every $\xi<\rho$ is therefore fixed, and evaluation gives
\[
 j(f(\xi))=j(f)(j(\xi))=f(\xi).
\]
Thus every element of $a$ is a fixed ordinal below $\lambda$. By
Lemma~\ref{lem:critical}, all of them are below $\kappa$, contradicting
unboundedness of $a$ in $\lambda$.
\end{proof}

\begin{lemma}[Relative exactingness excludes definable short cofinality]
\label{lem:relative}
If $\REx(\lambda,Y)$, no unbounded $a\subseteq\lambda$ of order type less than
$\lambda$ is ordinal definable from the set parameter $Y$.
\end{lemma}

\begin{proof}
Suppose such an $a$ exists. Use the canonical definable well-order of
$\OD_Y$ to select its least member $b$ that is an unbounded subset of
$\lambda$ of order type less than $\lambda$. Crucially, $b$ is now uniquely
definable from $Y$ and $\lambda$, without any additional free ordinal parameter.

This selection can be expressed by a single first-order formula. For example,
order definition codes first by a rank in which a unique definition holds,
then by the formula number and the finite ordinal-parameter tuple. The
property of being the first code defining an appropriate set is a set-theoretic
formula using satisfaction for \emph{set-sized} rank structures. No truth
predicate for $V$ is being assumed.

Choose a sufficiently high $V_\alpha$ reflecting that formula, its uniqueness,
and the relevant order-type and cofinality assertions. Also arrange that
$Y,b\in V_\alpha$ and that $\REx(\lambda,Y)$ supplies a witness there. Such
reflecting heights exist by the usual finite-formula Reflection Theorem.
In the expanded structure $(V_\alpha,Y)$, the set $Y$ is itself definable as
the set of all elements satisfying the predicate. Thus $Y\in X$ and
$j(Y)=Y$. Since $\lambda\in X$ is fixed, the unique definition of $b$ implies
$b\in X$ and $j(b)=b$. Lemma~\ref{lem:fixedcofinal} gives the contradiction.
\end{proof}

\begin{assessment}{Why the ordinal parameters cause no gap}
The original cofinal set may be defined using ordinals moved by a witness.
We never assume those ordinals are fixed. We first select a canonical least
\emph{cofinal set}, definable from the fixed pair $(Y,\lambda)$, and only then
choose a witness at a reflecting height. This is the specialized cofinality
version of the definability-transfer mechanism in the cover-exacting notes.
\end{assessment}

\section{Recovering an ultrafilter from a small trace}\label{sec:trace}

Here is the completeness calculation that ties the uniform cover bound to a
specific level of the HCD hierarchy. The underlying trace idea occurs in the
proof of the club-definability result in the Blue--Goldberg notes; we spell it
out at the exact endpoint required here. \cite[proof of Theorem 3.7]{cover}

\begin{lemma}[Endpoint trace coding]\label{lem:trace}
Let $\eta$ be an uncountable regular cardinal and let $U$ be an
$\eta$-complete ultrafilter on an ordinal $\theta$. For every sufficiently
large $\alpha$, there is a club $C_U\subseteq\Pow_\eta(V_\alpha)$ such that
\[
 \sigma\in C_U\quad\Longrightarrow\quad U\in\OD_\sigma.
\]
The parameter allowed on the right is the single set $\sigma$, together with
ordinal parameters.
\end{lemma}

\begin{proof}
Choose $\alpha$ large enough to contain $U$, all subsets of $\theta$, and the
few power-set and coding levels used below. Let $C_U$ be the club of
$\sigma\prec V_\alpha$ with $U,\theta\in\sigma$ and $|\sigma|<\eta$.

Fix $\sigma\in C_U$. Since $U\cap\sigma$ has fewer than $\eta$ members,
$\eta$-completeness gives
\[
 B_\sigma=\bigcap(U\cap\sigma)\in U.
\]
This set is nonempty because $U$ is proper. Define the ordinal
\[
 \beta=\min B_\sigma<\theta.
\]
For every $A\in\sigma$ with $A\subseteq\theta$, we have
\begin{equation}\label{eq:trace}
 A\in U\quad\Longleftrightarrow\quad\beta\in A.
\end{equation}
The forward direction is the definition of $\beta$. For the reverse direction,
if $A\notin U$, then $\theta\setminus A\in U$. Elementarity of $\sigma$ and
$A,\theta\in\sigma$ imply $\theta\setminus A\in\sigma$, so
$\beta\in\theta\setminus A$, contradicting $\beta\in A$.

We claim that $U$ is the unique ultrafilter $W$ on $\theta$ with $W\in\sigma$
satisfying
\begin{equation}\label{eq:decoder}
 \forall A\in\sigma\,
 [A\subseteq\theta\ \Longrightarrow\ (A\in W\leftrightarrow\beta\in A)].
\end{equation}
Existence follows from~\eqref{eq:trace}. For uniqueness, suppose $W\neq U$ is
another such ultrafilter in $\sigma$. In $V_\alpha$ there is a subset
$A\subseteq\theta$ on which $U$ and $W$ disagree. Because $U,W,\theta\in\sigma$
and $\sigma\prec V_\alpha$, such an $A$ can be chosen in $\sigma$.
Equation~\eqref{eq:decoder} would make both ultrafilters agree on $A$, a contradiction.

The uniqueness description uses only the set parameter $\sigma$ and the
ordinal parameters $\theta,\beta$. It is a definition in $V$; equivalently one
can use satisfaction for the expanded rank structure and include $\alpha$
as one more ordinal parameter. Therefore $U\in\OD_\sigma$.
\end{proof}

\begin{remark}[What is and is not being coded]\label{rem:trace}
The point $\beta$ need not belong to $\sigma$. The ultrafilter $U$ need not be
principal. Only its \emph{trace on subsets lying in $\sigma$} agrees with the
principal ultrafilter at $\beta$. The conclusion is not that $U$ itself is a
principal ultrafilter, or that its entire membership relation is determined
by one point without the parameter $\sigma$.
\end{remark}

\begin{corollary}[Finite parameters]\label{cor:finite}
If $a\in\CD(\eta)$ and $\eta$ is uncountable regular, then for every
sufficiently large $\alpha$ there is a club
$C\subseteq\Pow_\eta(V_\alpha)$ such that $a\in\OD_\sigma$ for all $\sigma\in C$.
\end{corollary}

\begin{proof}
Choose finitely many $\eta$-complete ultrafilters $U_0,\ldots,U_{m-1}$ and
ordinal parameters defining $a$. Choose one rank high enough for all of them,
and intersect the finitely many clubs supplied by Lemma~\ref{lem:trace}.
For each $\sigma$ in the intersection, substitute the definitions of all
$U_i$ from $\sigma$ into the definition of $a$. Only $\sigma$ and ordinal
parameters remain.
\end{proof}

Notice the exact inequality: for $\eta=\gamma^+$,
\[
 |\sigma|<\eta\quad\text{means}\quad |\sigma|\leq\gamma.
\]
Thus $\gamma^+$-completeness is enough. There is no need to move to
$\gamma^{++}$. Conversely, this calculation does not justify replacing
$\gamma^+$-completeness by $\gamma$-completeness when $|\sigma|=\gamma$.
This is an endpoint analysis of the proof, not a claim of a model-theoretically
optimal lower bound on completeness.

\section{The HCD regularity barrier}\label{sec:barrier}

\begin{theorem}[Complete-definability barrier]\label{thm:barrier}
Assume $\CEx_\gamma(\lambda)$, where $\gamma\geq\lambda$ is a cardinal.
If $a\subseteq\lambda$ is unbounded and $\ot(a)<\lambda$, then
\[
 a\notin\CD(\gamma^+).
\]
Consequently $\lambda$ is regular in $\HCD(\gamma^+)$.
\end{theorem}

\begin{proof}
Let $\eta=\gamma^+$ and suppose $a\in\CD(\eta)$ is as stated.
By Corollary~\ref{cor:finite}, at a sufficiently large rank there is a club
$C\subseteq\Pow_\eta(V_\alpha)$ on which $a$ is ordinal definable from
$\sigma$. By Fact~\ref{fact:stationary},
\[
 \{\sigma\in\Pow_\eta(V_\alpha):\REx(\lambda,\sigma)\}
\]
is stationary. Choose $\sigma$ in its intersection with $C$.
Then $a\in\OD_\sigma$ and $\REx(\lambda,\sigma)$, contradicting
Lemma~\ref{lem:relative}.

Let $N=\HCD(\eta)$. A cardinal of $V$ is a cardinal in the inner model $N$:
any inner-model bijection to a smaller ordinal would also be such a bijection
in $V$. If $\lambda$ were singular in $N$, it would have in $N$ a cofinal map
$f:\rho\to\lambda$ with $\rho<\lambda$. Its range $a$ belongs to $N$ and hence
to $\CD(\eta)$. Externally $|a|\leq|\rho|<\lambda$; since $a\subseteq\lambda$
and $\lambda$ is a cardinal, $\ot(a)<\lambda$. This contradicts the first part.
Thus $\cf^N(\lambda)=\lambda$.
\end{proof}

\begin{corollary}[Unbounded sets in the upper model are large]\label{cor:large}
Under the same hypotheses, every unbounded $b\subseteq\lambda$ with
$b\in\HCD(\gamma^+)$ has $\ot(b)=\lambda$ and $|b|^V=\lambda$.
The same statement holds in place of $\HCD(\gamma^+)$ for every inner model
contained in it.
\end{corollary}

\begin{proof}
The order type is at most $\lambda$. Theorem~\ref{thm:barrier} rules out every
smaller order type. Its increasing enumeration therefore gives a bijection
between $\lambda$ and $b$, in the inner model and in $V$.
For a smaller inner model, its sets are also sets of $\HCD(\gamma^+)$.
\end{proof}

\subsection{An equivalent forbidden-code formulation}

The first part of Theorem~\ref{thm:barrier} can be read without any inner-model
terminology. There cannot be a formula $\varphi$, a finite ordinal tuple
$\vec\xi$, finitely many $\gamma^+$-complete ultrafilters $\vec U$ on ordinals,
and a unique set $a$ such that
\[
 V\models\varphi(a,\vec U,\vec\xi),\qquad
 a\subseteq\lambda,\quad\sup a=\lambda,\quad\ot(a)<\lambda.
\]
The prohibition concerns the \emph{definability of the cofinal set}; the
ultrafilters themselves need not be elements of the exacting witness's domain.
The stationary collection of relative parameters is what bridges that gap.

\subsection{A necessary failure of small covering in the ambient universe}

Choose any increasing cofinal $\omega$-sequence in $\lambda$, and let $c$ be
its range. Such a sequence exists by Lemma~\ref{lem:critical}. If
$b\in\HCD(\gamma^+)$ contains $c$, then $b\cap\lambda$ is unbounded and belongs
to the same inner model. Corollary~\ref{cor:large} gives $|b|^V\geq\lambda$.
Thus the upper HCD model cannot even cover this countable set by one of its
sets of ambient size less than $\lambda$.

This observation alone does not refute cover exactingness. There is no theorem
in use here saying that $\HCD(\gamma^+)$ must have small covering in $V$.
The role of a smaller strongly compact cardinal is to put a short cofinal
\emph{set} in a different, lower-completeness HCD model, making the failure a
comparison between two canonical inner models.

\section{A smaller strongly compact cardinal forces a cofinality gap}
\label{sec:below}

\begin{theorem}[The local HCD gap]\label{thm:gap}
Suppose $\delta<\lambda\leq\gamma$, $\delta$ is strongly compact, and
$\CEx_\gamma(\lambda)$. Let
\[
 N_0=\HCD(\delta),\qquad N_1=\HCD(\gamma^+).
\]
There is an unbounded $a\subseteq\lambda$ such that
\begin{equation}\label{eq:a-gap}
 a\in N_0,\qquad \ot(a)<\delta,\qquad a\notin N_1.
\end{equation}
Moreover, for every $b\in N_1$ with $a\subseteq b$,
\begin{equation}\label{eq:no-small-cover}
 |b|^V\geq\lambda.
\end{equation}
In particular, $\cf^{N_0}(\lambda)<\delta$ but $\cf^{N_1}(\lambda)=\lambda$.
\end{theorem}

\begin{proof}
Let $c\subseteq\lambda$ be the range of an increasing cofinal $\omega$-sequence.
Since $\delta$ is uncountable and $N_0$ has the $\delta$-cover property by
Fact~\ref{fact:hcd}, choose a set of ordinals $a_0\in N_0$ containing $c$ with
$|a_0|^V<\delta$. Put $a=a_0\cap\lambda$. It belongs to $N_0$, is unbounded
in $\lambda$, and has ambient size less than $\delta$. As $\delta$ is a
cardinal, its increasing order type is less than $\delta$. This order type
and its increasing enumeration are absolute, so the enumeration also
witnesses $\cf^{N_0}(\lambda)<\delta$.

By Theorem~\ref{thm:barrier}, $N_1$ has no unbounded subset of $\lambda$ of
order type less than $\lambda$. Hence $a\notin N_1$. If $b\in N_1$ contains
$a$, then $b\cap\lambda$ is an unbounded $N_1$-subset of $\lambda$.
Corollary~\ref{cor:large} gives $|b\cap\lambda|^V=\lambda$, which implies
\eqref{eq:no-small-cover}. The remaining cofinality assertion is
Theorem~\ref{thm:barrier}.
\end{proof}

\subsection{A local covering principle that is already too much}

\begin{definition}[Generous local covering]\label{def:localcover}
For inner models $N_1\subseteq N_0$ and cardinals $\delta<\lambda$, let
$\Cov_\lambda^{\delta,<\lambda}(N_1,N_0)$ denote the assertion
\begin{equation}\label{eq:localcover}
 \begin{split}
 \forall a\in N_0\cap[\lambda]^{<\delta}\ \exists b\in N_1\quad
 a\subseteq b\subseteq\lambda\ \wedge\ |b|^V<\lambda.
 \end{split}
\end{equation}
Both sizes in this formula are measured in the ambient universe.
\end{definition}

This is much less demanding than agreement on all subsets of $\lambda$.
Even the permitted cover size may grow from less than $\delta$ to any cardinal
less than $\lambda$. The result does not depend on insisting on exact cardinal
preservation between the two inner models.

\begin{corollary}[Conditional inconsistency above a strongly compact cardinal]
\label{cor:localinconsistency}
The following theory is inconsistent:
\begin{equation}\label{eq:conditionaltheory}
 \begin{split}
 \ZFC+\exists\delta<\lambda\leq\gamma\ \bigl[
 &\SC(\delta)\wedge\CEx_\gamma(\lambda)\\[-2pt]
 &{}\wedge\Cov_\lambda^{\delta,<\lambda}
       (\HCD(\gamma^+),\HCD(\delta))\bigr].
 \end{split}
\end{equation}
The same conclusion holds with the covering clause replaced by
\begin{equation}\label{eq:localagreement}
 \HCD(\delta)\cap[\lambda]^{<\delta}\subseteq\HCD(\gamma^+).
\end{equation}
\end{corollary}

\begin{proof}
Apply~\eqref{eq:localcover} to the set $a$ from Theorem~\ref{thm:gap}. The
resulting $b$ contradicts~\eqref{eq:no-small-cover}. Condition
\eqref{eq:localagreement} implies~\eqref{eq:localcover} by choosing $b=a$.
\end{proof}

Ordinary $\delta$-covering of $N_0$ by $N_1$, with covers of $N_1$-size less
than $\delta$, also implies~\eqref{eq:localcover} for the relevant small sets
whenever it applies to them. More precisely, it suffices that it apply to
sets that have $N_0$-size less than $\delta$; the particular $a$ in the proof
has such size because its order type is less than $\delta$ in $N_0$.
Thus the theorem rules out the standard internal version as well, without
assuming that all ambient-small $N_0$-sets are internally small.

\subsection{Recovering the extendible prohibition}

\begin{corollary}[Known boundary recovered]\label{cor:extendible}
An extendible $\delta$ cannot lie below a $\gamma$-cover-exacting $\lambda$.
\end{corollary}

\begin{proof}
An extendible cardinal is strongly compact. By Fact~\ref{fact:hcd},
$\HCD(\delta)=\HCD$. Since $\delta<\gamma^+$, monotonicity gives
\[
 \HCD(\delta)=\HCD\subseteq\HCD(\gamma^+)
 \subseteq\HCD(\delta).
\]
Hence the two models are equal and~\eqref{eq:localagreement} holds, contrary
to Corollary~\ref{cor:localinconsistency}.
\end{proof}

This is a recovery of Blue--Goldberg's Theorem 3.6, not a second claim to that
result. \cite{cover} The new formulation separates what is used from strong
compactness---a short cofinal set in $\HCD(\delta)$---from what extendibility
adds---stabilization of the complete-definability hierarchy.

\subsection{Exactly what is missing in the unrestricted problem}

Strong compactness by itself supplies $\delta$-covering from $V$ into
$\HCD(\delta)$. It does \emph{not}, by the theorem cited here, supply covering
from $\HCD(\delta)$ into $\HCD(\gamma^+)$. These are different pairs of models.
Using the former theorem as though it were the latter would erase the essential
gap in~\eqref{eq:cofgap}.

A resolution by this route would require proving a principle such as
\eqref{eq:localcover} from the original coexistence hypotheses, or obtaining
one forbidden short cofinal code in $\CD(\gamma^+)$ by another method.
The present proof does neither. Its output is the explicit incompatibility
\eqref{eq:conditionaltheory} and the necessary, sharply witnessed failure of
covering in every putative model of the unrestricted coexistence theory.

\section{A larger strongly compact cardinal forces a proper ground}
\label{sec:ground}

The same regularity theorem gives an incompatibility with a familiar axiom
that does not mention HCD. Here the strongly compact cardinal lies above the
\emph{cover bound}, not below $\lambda$.

\begin{theorem}[Ground singularization]\label{thm:ground}
Suppose $\lambda\leq\gamma<\kappa$, $\CEx_\gamma(\lambda)$, and
$\kappa$ is strongly compact. Then
\[
 N=\HCD(\kappa)
\]
is a proper set-forcing ground of $V$, and
\begin{equation}\label{eq:groundcof}
 \cf^N(\lambda)=\lambda,\qquad \cf^V(\lambda)=\omega.
\end{equation}
For every presentation $V=N[G]$ with forcing $\mathbb P\in N$,
\begin{equation}\label{eq:forcinglower}
 |\mathbb P|^N\geq\lambda.
\end{equation}
\end{theorem}

\begin{proof}
Because $\kappa$ is a cardinal greater than $\gamma$, it is at least
$\gamma^+$. By monotonicity,
\[
 N=\HCD(\kappa)\subseteq\HCD(\gamma^+).
\]
Fact~\ref{fact:hcd} says that $N\models\ZFC$ and that $N$ is a set-forcing
ground of $V$. A cofinal map in $N$ from an ordinal less than $\lambda$ into
$\lambda$ would belong to $\HCD(\gamma^+)$ and contradict
Theorem~\ref{thm:barrier}. Thus $\lambda$ is regular in $N$.
Lemma~\ref{lem:critical} gives its ambient cofinality $\omega$, so $N\neq V$.

For the size assertion, suppose $\mu=|\mathbb P|^N<\lambda$.
In $N$, choose a name $\dot f$ and a condition $p\in G$ forcing that
$\dot f:\omega\to\lambda$ is cofinal. For each $n<\omega$, choose a maximal
antichain below $p$ deciding $\dot f(n)$. Each antichain has at most $\mu$
elements, and hence there are at most $\mu$ possible values for each $n$.
Let $B\in N$ be the union of these sets of possible values. Inside $N$,
\[
 |B|\leq\mu\cdot\aleph_0<\lambda.
\]
Regularity of $\lambda$ in $N$ implies $\sup B<\lambda$. But $p$ forces
$\ran(\dot f)\subseteq B$, contradicting cofinality. This proves
\eqref{eq:forcinglower}.
\end{proof}

\begin{corollary}[Ground Axiom incompatibility]\label{cor:ga}
The theory
\begin{equation}\label{eq:ga}
 \ZFC+\GA+
 \exists\lambda\leq\gamma<\kappa\,
 [\CEx_\gamma(\lambda)\wedge\SC(\kappa)]
\end{equation}
is inconsistent. In particular, the Ground Axiom together with a proper class
of strongly compact cardinals excludes every cover-exacting cardinal with a
set-sized uniform cover bound.
\end{corollary}

\begin{proof}
Theorem~\ref{thm:ground} supplies a proper ground, contradicting $\GA$.
With a proper class of strongly compact cardinals, choose one above any
given bound $\gamma$.
\end{proof}

This is a conditional incompatibility among explicitly named principles.
It does not imply that cover exactingness and a larger strongly compact
cardinal are inconsistent: it implies that their coexistence requires a
nontrivial ground. The ground need not be HOD, and the forcing witnessing it
need not be Prikry forcing. The proof only determines its effect on
$\cf(\lambda)$ and the lower bound~\eqref{eq:forcinglower}.

There is also a useful consistency check on the ordering. The cover-exacting
notes describe constructions with cover-exacting cardinals below larger
large cardinals. \cite[discussion following Theorem 3.5]{cover} Nothing here
contradicts those constructions; a nontrivial forcing extension is precisely
the sort of universe the theorem permits.

\section{An I0-equiconsistent sharp HOD boundary}\label{sec:positive}

The preceding theorems are negative only when a specified covering or geology
assumption is added. We now make the positive side equally explicit. This
section uses an existing forcing construction and derives a stronger package
of boundary properties from its output. It does not claim a new construction
of an ultraexacting cardinal.

\subsection{The two consistency inputs}

Aguilera--Bagaria--Goldberg--L\"ucke prove both
\begin{equation}\label{eq:knowni0}
 \Con(\ZFC+\exists\lambda\,\UEx(\lambda))
 \longleftrightarrow\Con(\ZFC+\Izero)
\end{equation}
and the following preservation-and-coding result: from an ultraexacting
$\lambda$, forcing can preserve its ultraexactingness and $V_\lambda$ while
making $V_\lambda\subseteq\HOD$. \cite[Theorem 4.5; Proposition 3.9]{abgl}
We also use the established theorem that an exacting $\lambda$ is regular
in HOD. \cite[Theorem 2.10]{abl}

The consistency transformations in~\eqref{eq:knowni0} need not preserve the
ordinal at which the principle holds. Our equiconsistency conclusion below
does not assert otherwise.

\subsection{The boundary theory}

\begin{definition}\label{def:boundary}
Let $\Bdry$ assert that there are a cardinal $\lambda$ and an increasing
cofinal $\omega$-sequence with range $c\subseteq\lambda$ satisfying all of the
following:
\begin{enumerate}
\item $\lambda$ is ultraexacting and there is no exacting cardinal below $\lambda$.
\item $V_\lambda\subseteq\HOD$, but $c\notin\HOD$. In particular,
      $V_{\lambda+1}\nsubseteq\HOD$.
\item $\lambda$ is regular in HOD, and for every $a\in\HOD$ with
      $|a|^{\HOD}<\lambda$, the intersection $c\cap a$ is finite and belongs to HOD.
\item Every $b\in\HOD$ with $c\subseteq b$ has $|b|^V\geq\lambda$.
\end{enumerate}
\end{definition}

All cardinalities explicitly superscripted by HOD are computed internally.
The final cardinality is an ambient one. These two conventions are not
interchangeable in a general inner-model comparison.

\begin{lemma}[A single cofinal set supplies all the boundary witnesses]
\label{lem:boundary}
Suppose $\lambda$ is exacting and $V_\lambda\subseteq\HOD$. Then any range $c$
of an increasing cofinal $\omega$-sequence in $\lambda$ satisfies items
(2)--(4) of Definition~\ref{def:boundary}, and there is no exacting cardinal
below $\lambda$.
\end{lemma}

\begin{proof}
Let $H=\HOD$. The cited exacting regularity theorem gives
$\cf^H(\lambda)=\lambda$, whereas Lemma~\ref{lem:critical} gives
$\cf^V(\lambda)=\omega$.
If $c\in H$, its increasing enumeration would be a cofinal map
$\omega\to\lambda$ in $H$, a contradiction. Since $c$ is unbounded in
$\lambda$, its rank is exactly $\lambda$. Thus it belongs to
$V_{\lambda+1}\setminus H$, while every set of smaller rank is in $H$ by
assumption. This proves item (2) and the regularity part of (3).

Now let $a\in H$ with $|a|^H<\lambda$. The set $a\cap\lambda$ is in $H$ and
has $H$-size less than $\lambda$. Regularity in $H$ makes it bounded in
$\lambda$. An increasing $\omega$-sequence cofinal in $\lambda$ has only
finitely many terms below any particular bound. Hence $c\cap a$ is finite.
Every finite set of ordinals is in HOD, using its finitely many members as
ordinal parameters. This proves (3).

If $b\in H$ contains $c$, the set $d=b\cap\lambda$ is an unbounded
$H$-subset of $\lambda$. Its order type cannot be less than $\lambda$, by
regularity in $H$, and cannot exceed $\lambda$, since $d\subseteq\lambda$.
Thus $\ot(d)=\lambda$. The increasing enumeration is an actual bijection
between $\lambda$ and $d$, so $|b|^V\geq\lambda$, proving (4).

Finally, suppose $\mu<\lambda$ were exacting. It would have ambient
cofinality $\omega$ and be regular in HOD. The range of an increasing cofinal
$\omega$-sequence in $\mu$ has rank $\mu<\lambda$, so it belongs to
$V_\lambda\subseteq H$. Its increasing enumeration in $H$ contradicts the
regularity of $\mu$ there. Therefore no such $\mu$ exists.
\end{proof}

\begin{theorem}[Equiconsistency of the full boundary package]\label{thm:eq}
\[
 \Con(\ZFC+\Bdry)\quad\longleftrightarrow\quad\Con(\ZFC+\Izero).
\]
\end{theorem}

\begin{proof}
For the forward implication, $\Bdry$ includes the existence of an ultraexacting
cardinal. Apply the corresponding direction of~\eqref{eq:knowni0}.

Conversely, the other direction of~\eqref{eq:knowni0} provides the relative
consistency of a model of $\ZFC$ with an ultraexacting $\lambda$.
Apply the cited preservation-and-coding forcing to obtain a model in which
$\lambda$ remains ultraexacting and $V_\lambda\subseteq\HOD$.
Choose an increasing cofinal $\omega$-sequence in $\lambda$ in that model.
Lemma~\ref{lem:boundary} verifies every additional clause of $\Bdry$.
Thus the original consistency assumption gives $\Con(\ZFC+\Bdry)$.
\end{proof}

The first-exacting conclusion and the optimal rank agreement are already
present in the cited paper's discussion and Corollary 3.10. \cite{abgl}
The role of the formulation above is to include a single, explicit witness
to all the approximation and covering failures, and to record the
consistency strength of the whole package. It should not be mistaken for an
independent proof of~\eqref{eq:knowni0}.

\subsection{Approximation failure, with the quantifiers visible}

For inner models $N\subseteq W$, the $\rho$-approximation property says that
whenever $A\in W$, $A\subseteq N$, and
\[
 A\cap a\in N\quad\text{for every }a\in N\text{ with }|a|^N<\rho,
\]
then $A\in N$. It is customary to study this for regular uncountable
$\rho$, but the displayed condition makes sense at any cardinal parameter.

In a model of $\Bdry$, the same set $c\subseteq\lambda\subseteq\HOD$ satisfies
\begin{equation}\label{eq:apxwitness}
 c\notin\HOD,\qquad
 \forall a\in\HOD\,
 [|a|^{\HOD}<\lambda\ \Longrightarrow\ c\cap a\in\HOD].
\end{equation}
It follows that $(\HOD,V)$ fails the $\rho$-approximation property for every
ambient regular uncountable $\rho<\lambda$. Equation~\eqref{eq:apxwitness}
also witnesses failure of the displayed approximation condition at $\lambda$
itself, without calling the ambiently singular $\lambda$ regular.

Likewise, item (4) prevents HOD from covering the countable set $c$ by any
of its sets of ambient size less than $\lambda$. In particular, the usual
$\rho$-cover property fails for every ambient regular uncountable
$\rho<\lambda$. An internally small cover would also be externally small,
so the internal-size convention cannot avoid this conclusion.

\subsection{What the positive package does not imply}

It does not put a short cofinal map into HOD. It does not imply agreement on
subsets of $V_\lambda$. It does not make an arbitrary ambient well-order of
$V_{\lambda+1}$ ordinal definable. And it does not turn an ultraexacting
cardinal into a cover-exacting cardinal with any particular cover bound.
The affirmative construction therefore neither contradicts
Theorem~\ref{thm:barrier} nor solves the smaller-strongly-compact coexistence
question.

Its exact boundary is useful: every set of rank \emph{less than} $\lambda$
may be in HOD, but one cannot extend that inclusion to all sets of rank
$\lambda$. The countable cofinal range already supplies a counterexample.

\section{What has been established, and what remains open}\label{sec:status}

\subsection{The contribution relative to the supplied report}

The supplied report proposed locating a local definability condition behind
the extendible obstruction. The present work gives such a condition with
explicit cardinal parameters. It proves regularity at the particular HCD
level $\gamma^+$, extracts a short cofinal set from the smaller strongly compact
cardinal's HCD model, and rules out even generous local covering between those
two levels. The opposite ordering produces a further named-axiom
incompatibility with the Ground Axiom.

The core statements are
\[
 \begin{gathered}
 \CEx_\gamma(\lambda)\Rightarrow
       \cf^{\HCD(\gamma^+)}(\lambda)=\lambda,\\
 \SC(\delta)\wedge\delta<\lambda\wedge\CEx_\gamma(\lambda)
       \Rightarrow\neg\Cov_\lambda^{\delta,<\lambda}
        (\HCD(\gamma^+),\HCD(\delta)),\\
 \SC(\kappa)\wedge\gamma<\kappa\wedge\CEx_\gamma(\lambda)
       \Rightarrow\neg\GA.
 \end{gathered}
\]
The proofs are given above, with imported results isolated in
Section~\ref{sec:setup}. These are mathematical consequences, not proposals
for future lemmas. Their dependence on the lecture-note stationary
characterization is explicit.

\subsection{Novelty and verification status}

The bounded HCD regularity, local covering obstruction, and ground
singularization statements are derived in this follow-up from established
techniques. I did not locate these exact formulations in the primary texts
checked. That is not a certification that no equivalent result is known;
some or all may already be recognized corollaries by specialists. The
ultrafilter trace argument is credited to the notes, and the complete
definability, covering, and ground theorems are credited to Goldberg.

The $\Izero$ theorem here is deliberately described as an equiconsistency
\emph{refinement/package}: the nontrivial consistency construction is already
in Aguilera--Bagaria--Goldberg--L\"ucke. No historical novelty is claimed for
the forcing construction, the first-exacting corollary, or the impossibility
of extending $V_\lambda\subseteq\HOD$ to $V_{\lambda+1}$.

The arguments have been checked at the level of ordinary mathematical proofs,
including the completeness endpoint, ordinal parameters, internal versus
external size, and the direction of covering. They have not been independently
refereed or formalized in a proof assistant. Typesetting compilation is not
presented as mathematical machine verification.

\subsection{The remaining mathematical obstruction}

The theory
\[
 \ZFC+\exists\delta<\lambda\leq\gamma\,
       [\SC(\delta)\wedge\CEx_\gamma(\lambda)]
\]
is not refuted by this paper. Under that theory, the HCD tower must have the
specific cofinality gap in Theorem~\ref{thm:gap}. To obtain an unconditional
refutation by this route, one must show that this gap is impossible using only
the original hypotheses. The published covering theorem does not do that:
it covers into the lower-completeness model, whereas the contradiction needs
a cover in the higher-completeness model.

Conversely, a positive model of the theory would necessarily contain the
short set $a$ in~\eqref{eq:a-gap} and exhibit its failure to acquire any small
cover at the upper HCD level. This is a concrete requirement on a construction,
not just the vague requirement that HOD be far from $V$.

The global cardinal-preserving embedding route was also checked against the
primary literature. Reconstructing the embedding range from the external
stationarity pattern of a transported partition is already part of
Hamkins--Kirmayer--Perlmutter's argument. \cite[proofs of Theorems 11--12]{hkp}
It is therefore not offered here as a new result. No new inconsistency of
$\Izero^{\#}$, ordinary exacting--extendible coexistence, or HOD-Berkeley
cardinals is asserted.

\begin{assessment}{Bottom line}
There is a proved local inconsistency theorem, a proved Ground Axiom
incompatibility, and a fully specified equiconsistency corollary. The main
newly articulated boundary is the jump from short cofinality in
$\HCD(\delta)$ to regularity in $\HCD(\gamma^+)$, with an actual set witnessing
the failure of every cover of size less than $\lambda$. Removing that
covering failure from a putative coexistence model would require additional
mathematics; it is not supplied by changing the name of the compactness axiom.
\end{assessment}

\clearpage
\appendix
\section{Dependency and hypothesis audit}\label{app:audit}

\begingroup\small
\renewcommand{\arraystretch}{1.13}
\begin{longtable}{@{}P{.24\textwidth}P{.34\textwidth}P{.35\textwidth}@{}}
\toprule
\textbf{Conclusion}&\textbf{Exact additional input}&\textbf{Potential invalid substitution avoided}\\
\midrule\endfirsthead
\toprule
\textbf{Conclusion}&\textbf{Exact additional input}&\textbf{Potential invalid substitution avoided}\\
\midrule\endhead
Critical sequence cofinal in $\lambda$ & Local Kunen theorem in ambient ZFC; $\lambda$ a cardinal. & A rank-$\lambda+1$ embedding is not refuted. The forbidden restriction occurs only if the critical limit lies strictly below the cardinal height.\\[6pt]
Relative definable cofinality contradiction & Witnesses at arbitrarily high reflecting ranks; canonical least cofinal set. & Arbitrary ordinal parameters of an original definition are not presumed fixed.\\[6pt]
Trace coding & $\eta$ regular; $|\sigma|<\eta$; $U$ $\eta$-complete; $\sigma\prec V_\alpha$. & Principal behavior on a small trace is not principality of $U$. The coding point need not lie in $\sigma$.\\[6pt]
Regularity in $\HCD(\gamma^+)$ & Stationary relative-witness characterization, not merely one cover. & Internal stationarity is not used as ambient stationarity. No Choice in this HCD level is assumed.\\[6pt]
Short cofinal set in $\HCD(\delta)$ & $\delta<\lambda$ strongly compact; the published $\delta$-cover theorem. & Covering into $\HCD(\delta)$ is not covering into $\HCD(\gamma^+)$.\\[6pt]
Local inconsistency & Local covering or local small-set agreement between the two specified HCD levels. & That extra principle is not inferred from strong compactness.\\[6pt]
Extendible contradiction & Published $\HCD(\delta)=\HCD$ for extendible $\delta$. & This recovers a known theorem; it is not a new refutation of ordinary exactingness.\\[6pt]
Proper-ground theorem & A strongly compact $\kappa>\gamma$, with $\gamma$ the fixed cover bound. & The order $\delta<\lambda$ from the other theorem is not reused. The ground is HCD, not necessarily HOD.\\[6pt]
Equiconsistency package & Published ultraexacting/$\Izero$ comparison and high coding that preserves ultraexactingness. & The same ordinal need not witness $\Izero$ in the consistency transformation. No new forcing proof is claimed.\\[6pt]
Approximation witness & HOD regularity at $\lambda$, a strictly increasing cofinal $\omega$-sequence, and internal smallness of test sets. & Ambiently singular $\lambda$ is not called regular; the witness formula is stated directly.\\
\bottomrule
\end{longtable}
\endgroup

\subsection{The core dependency chain}

The proof of the principal negative result can be checked in the following order:
\[
 \begin{gathered}
 \text{local Kunen}\ \Longrightarrow\ \text{Lemma~\ref{lem:critical}}
 \ \Longrightarrow\ \text{Lemmas~\ref{lem:fixedcofinal}--\ref{lem:relative}},\\[3pt]
 \text{ultrafilter completeness}\ \Longrightarrow\ \text{Lemma~\ref{lem:trace}}
 \ \Longrightarrow\ \text{Corollary~\ref{cor:finite}},\\[3pt]
 \text{Fact~\ref{fact:stationary}}+
 \text{Lemma~\ref{lem:relative}}+\text{Corollary~\ref{cor:finite}}
 \ \Longrightarrow\ \text{Theorem~\ref{thm:barrier}},\\[3pt]
 \text{Theorem~\ref{thm:barrier}}+\text{strongly compact HCD covering}
 \ \Longrightarrow\ \text{Theorem~\ref{thm:gap}},\\[3pt]
 \text{Theorem~\ref{thm:gap}}+\text{local covering}
 \ \Longrightarrow\ \bot.
 \end{gathered}
\]
The geology theorem replaces the last two lines by the published HCD ground
theorem at a strongly compact cardinal above $\gamma$. The affirmative
$\Izero$ package is independent of this dependency chain.

\section{Source and reproducibility notes}\label{app:sources}

The starting document was the supplied
\nolinkurl{Large_Cardinals_Unified_Report.tex}. Its preamble was reused for
the body and mathematical fonts, page geometry, heading hierarchy, color
values, running furniture, and callout style. The new document does not
reproduce the earlier candidate survey.

The primary sources were checked on 18 September 2026. In particular,
Goldberg's published HCD theorem was checked directly rather than relying
on the stronger hypothesis in its short lecture-note presentation.
The cover-exacting stationary characterization remains attributed to the
July 2026 lecture notes. The ABGL results are cited to arXiv version 1; the
ABL regularity theorem is cited to version 4.

The archive contains this source, its compiled PDF, and a README with build
instructions and the verification scope. The source is self-contained apart
from standard LaTeX packages. Compile twice with pdfLaTeX, or use
\texttt{latexmk -pdf}. The required body/math packages are \texttt{newpxtext}
and \texttt{newpxmath}; the sans-serif accents follow the supplied preamble.
No external figures, bibliography database, or custom font files are required.

\vspace{10pt}
\phantomsection
\addcontentsline{toc}{section}{References}
\begin{thebibliography}{99}
\small
\setlength{\itemsep}{7pt}
\setlength{\parskip}{1pt}

\bibitem{previous}
\emph{Large cardinals at the inconsistency frontier: A unified assessment of
ZFC-compatible hypotheses, their obstructions, and the evidence on both sides}.
Report supplied by the user, dated 18 September 2026.
Source: \nolinkurl{Large_Cardinals_Unified_Report.tex}.
Sections 6 and 13 give the cover-exacting target and proposed theorem-transfer programme.

\bibitem{cover}
G.~Goldberg. \emph{Consistency beyond the Kunen inconsistency}.
Lecture notes dated 1 July 2026, reporting joint work with D.~Blue.
\url{https://math.berkeley.edu/~goldberg/Slides/CoverExact.pdf}.
Definitions 2.6, 3.1, and 3.2; Lemma 3.3; Proposition 3.4;
Theorems 3.5--3.8; Problem 5.2.
The stationary characterization and trace-coding mechanism are the inputs
used here. Cited as lecture notes, not as a refereed cover-exacting paper.

\bibitem{unique}
G.~Goldberg. \emph{The uniqueness of elementary embeddings}.
Journal of Symbolic Logic \textbf{89}(4) (2024), 1430--1454.
Author manuscript:
\url{https://math.berkeley.edu/~goldberg/Papers/UniquenessOfElementaryEmbeddings.pdf}.
Propositions 4.2--4.3; Theorems 4.4, 4.10, 4.14, and 4.15.
The ground and covering assertions require strong compactness; the
stabilization assertion uses extendibility.

\bibitem{abl}
J.~P.~Aguilera, J.~Bagaria, and P.~L\"ucke.
\emph{Large cardinals, structural reflection, and the HOD Conjecture}.
\arxiv{2411.11568v4}, 16 September 2025.
Theorem 2.10 supplies regularity in $\HOD_{V_\lambda}$, and hence in HOD,
for an exacting cardinal $\lambda$.

\bibitem{abgl}
J.~P.~Aguilera, J.~Bagaria, G.~Goldberg, and P.~L\"ucke.
\emph{Large cardinals beyond HOD}.
\arxiv{2509.10254v1}, 12 September 2025.
Proposition 3.9, Corollary 3.10, and Theorem 4.5 supply the high coding,
first-exacting control, and ultraexacting/$\Izero$ equiconsistency used in
Section~\ref{sec:positive}.

\bibitem{kunen}
K.~Kunen. \emph{Elementary embeddings and infinitary combinatorics}.
Journal of Symbolic Logic \textbf{36}(3) (1971), 407--413.
\doi{10.2307/2269948}.
The local rank-$\nu+2$ obstruction in the critical-sequence argument is the
classical inconsistency input.

\bibitem{hkp}
J.~D.~Hamkins, G.~Kirmayer, and N.~L.~Perlmutter.
\emph{Generalizations of the Kunen inconsistency}.
Annals of Pure and Applied Logic \textbf{163} (2012), 1872--1890.
\arxiv{1106.1951v2}.
Used for the scope of the classical obstruction and to identify the
stationarity-pattern reconstruction as an existing argument rather than
an additional new result.

\end{thebibliography}
\end{document}
